\documentclass[12pt, aspectratio=169]{beamer}

\usetheme{Madrid}
\usecolortheme{default}

\usepackage[utf8]{inputenc}
\usepackage[T1]{fontenc}
\usepackage[brazil]{babel}
\usepackage{tcolorbox}


\title{Sobre \textit{Overfitting} e Generalização de Modelos}
\author{J. Roberto S. de Moura}
\date{\today}


\begin{document}

\begin{frame}
  \titlepage
\end{frame}


\begin{frame}\frametitle{Agenda}
  \tableofcontents
\end{frame}

\section{Motivação}
\begin{frame}\frametitle{Motivação}
  Imagine um estudante que decora todas as provas anteriores de um professor,
  mas que não aprende o conteúdo.\\[1cm]
\centering\fbox{
    \textbf{O que acontece com ele na próxima prova?}
}
\end{frame}

\section{Propósito de um modelo de \textit{Machine Learning}}
\begin{frame}\frametitle{Propósito de um modelo de \textit{Machine Learning}}

  \begin{tcolorbox}[colback=gray!10, colframe=black!70, title=Problema Central]
  
    Um modelo de ML minimiza uma função de perda nos dados de treino.
    Esse é o único objetivo dele. Ele não sabe o que é "generalizar".
    Ele não sabe que existe um mundo além dos dados que você deu.

      \begin{enumerate}
        \item O que realmente queremos: minimizar o erro em dados \textbf{NOVOS} - nunca vistos
        \item O que o modelo está tentando fazer: minimizar o erro nos dados de \textbf{TREINO}
      \end{enumerate}

  \end{tcolorbox}

\end{frame}

\section{O que é o \textit{Overfitting}?}
\begin{frame}\frametitle{O que é o \textit{Overfitting}?}
Quando passamos os dados para o modelo, ele não sabe o que é o ruído que está dentro dos dados.
$$
\Downarrow
$$
\textbf{Ele aprende o ruído dos dados de treino como se fosse sinal.}

Vale lembrar que os dados reais, possuem duas componentes:
$$
  y = f(x) + \epsilon
$$
  sendo $f(x)$ o padrão real que queremos aprender e $\epsilon$ um ruído aleatório.

\end{frame}

\section{O que é a Generalização?}
\begin{frame}\frametitle{O que é a Generalização?}
Generalizar é a capacidade de um modelo de ML de performar bem em dados novos, nunca vistos antes.\\[1cm]
\textbf{O que é um modelo que generaliza?}
\begin{tcolorbox}[colback=gray!10, colframe=black!70, title=Definição de Generalização]
  Um modelo que generaliza é aquele que consegue aprender o padrão real $f(x)$, e não o ruído $\epsilon$.
\end{tcolorbox}
\end{frame}

\section{Um Pequeno Experimento Mental}
\begin{frame}\frametitle{Um Pequeno Experimento Mental}
Situação: Dois médicos diagnosticando pneumonia.\\[0.5cm]
\begin{table}
  \begin{tabular}{l|l}

    Médico A & Médico B \\
    \hline
    Estudou 10.000 casos variados & Memorizou 100 casos em detalhes absurdos \\
    Aprende padrões gerais & Aprende ``paciente José, 45 anos, fuma'' \\
    Erra um pouco em casos conhecidos & Acerta todos os casos conhecidos \\
    Funciona com pacientes novos & Falha miseravelmente com pacientes novos \\

  \end{tabular}
\end{table}
Conclusão: O médico A é um modelo que generaliza, enquanto o médico B é um modelo que sofre de \textit{overfitting}
\end{frame}

\begin{frame}\frametitle{Mas Por Que o \textit{Overfitting} Acontece?}
Três causas estruturais:
\begin{itemize}
  \item[$\rightarrow$] Capacidade excessiva do modelo: Muitos parâmetros para pouco dado. O modelo tem "espaço" para memorizar.
  \item[$\rightarrow$] Dados insuficientes ou não representativos: Pouco dado = o modelo não vê variação suficiente para aprender o padrão real.
  \item[$\rightarrow$] Treinamento longo demais: Nas primeiras épocas, o modelo aprende padrões gerais. Depois, começa a ajustar detalhes específicos dos dados de treino.\\
\begin{itemize}
      \item[$\rightarrow$] época 10:  $val\_loss$ cai   $\leftarrow$ está generalizando
      \item[$\rightarrow$] época 50:  $val\_loss$ sobe  $\leftarrow$ está memorizando
\end{itemize}
\end{itemize}
\end{frame}



\end{document}
'